\section{Discussion}
\label{sec:discussion}
\todo{discussion}

\subsection{Limitations}
% 27 runs; one prompt design; models change; the Latin-square assumption of no
% seat x rotation interaction; judges are models too; cost-driven run length.
%
% Recognition of the games (Germ, 1 Nov: keep the design, report it).
% Finding: the recognition check (zhou2025pimmur) before the freeze. All 15
% answers name the prisoner's dilemma and a principal-agent or corruption
% model; 14 name the inspection game; all guess the study's aims. All 5 Claude
% play answers derive crime 4/15, inspection 1/4 and b > 100a + 11 from the
% stated payoffs; most Gemini and GPT answers derive the first two.
% Risk: a result on the theory line is what recall would produce too
% (barrie2025leakage), so a match is not evidence of reasoning.
% Mitigations to state:
% - The prompts name no game; one city framing for every seat and run (lore2024).
% - H6 compares models and seats; recognition is at ceiling in all three
%   models, so it cannot explain a model difference (flintashery2025reply:
%   model-dependent dynamics stay open to study under leakage).
% - H1 to H4 measure departures from theory; recognition did not produce
%   equilibrium plans (Claude's Chief plans full inspection on day 1).
% - The equilibrium numbers are Tidewall's own (4/15, 1/4, b > 100a + 11),
%   not textbook values; recall gives the method, not the number.
% - Recognition rate and the diary scan for game-theory terms reported by
%   model and seat; matched bot baselines on each seed.
% Not done, named as future work: counterfactual payoffs and labels
% (georgousis2026counterfactual); disrupting the history's statistical
% structure against a history-free benchmark (liu2026recursive).
\todo{limitations}
